\chapter{Introduction}

\section{Background}

Policy decisions affecting millions of citizens are routinely made with limited understanding of how diverse stakeholders will respond. Environmental impact assessments in the United States require a median of 2.2 years to complete~\cite{ceq2024}, while professional stakeholder war-gaming exercises conducted by firms such as McKinsey command fees exceeding \$500,000 per engagement. These constraints---temporal, financial, and logistical---severely limit the ability of policymakers to anticipate opposition, identify coalition opportunities, and stress-test policy proposals before implementation.

The advent of Large Language Models (LLMs) has opened a fundamentally new approach to this problem. Rather than assembling physical participants or relying on rule-based simulations, LLMs can serve as synthetic stakeholder agents capable of naturalistic dialogue, nuanced argumentation, and context-sensitive negotiation. Recent demonstrations have shown that LLM-based agents can replicate individual human responses with up to 85\% accuracy on attitudinal surveys~\cite{park2024generative}, and that AI-generated consensus statements are preferred over human-mediated ones by 56\% of participants~\cite{bakker2024habermas}. The commercial viability of this approach is evidenced by substantial venture funding: Simile, a Stanford spin-off, raised \$100 million in Series~A funding for LLM-based persona simulation, while Aaru achieved a valuation approaching \$1 billion for synthetic population prediction. Gartner has classified ``Decision Intelligence'' as a transformational technology in its 2025 AI Hype Cycle.

Despite this momentum, no production-grade tool exists for simulating how a specific group of 5--10 stakeholders would deliberate around a policy table---what arguments they would advance, where they would resist, and how compromise might emerge. This gap is not due to insufficient market demand but rather to unresolved technical challenges that render current LLM-based simulations unreliable for substantive policy analysis.


\section{Problem Statement}

This study identifies three interconnected problems that prevent LLM-based multi-agent simulation from producing meaningful policy deliberation.

\textbf{Problem 1: Persona Drift.} When an LLM is assigned a specific personality profile---for example, a fiscally conservative regulator with high Conscientiousness and low Openness---it tends to deviate from that profile as dialogue progresses. Over multiple turns of interaction, the agent's behavioral characteristics converge toward a generic, agreeable mean, eroding the distinctive perspectives that make stakeholder simulation informative. Without mechanisms to monitor and correct this drift, the diversity of viewpoints that motivates multi-stakeholder simulation is lost within the first few turns of dialogue.

\textbf{Problem 2: Sycophancy Convergence.} LLMs trained through RLHF are optimized to produce responses that human evaluators rate favorably, which creates an inherent bias toward agreement and validation~\cite{wei2024sycophancy}. In a multi-agent deliberation context, this manifests as rapid consensus formation: agents default to acknowledging each other's points, expressing agreement, and converging on shared positions regardless of their assigned stakeholder roles. Jain~\etal~\cite{jain2026sycophancy} demonstrate that interaction context further amplifies sycophantic tendencies in LLMs. An unconstrained simulation of a labor policy debate, for instance, may reach unanimous agreement within three turns---an outcome that is neither realistic nor informative.

\textbf{Problem 3: RLHF Behavioral Ceiling.} RLHF safety training structurally suppresses certain personality characteristics. Specifically, traits associated with emotional volatility (high Neuroticism) and confrontational behavior (low Agreeableness) are difficult or impossible for LLMs to express, as the safety alignment process penalizes outputs containing anxiety, self-doubt, or aggressive disagreement. This creates a fundamental ceiling on the range of stakeholder personas that can be faithfully simulated, regardless of the sophistication of the surrounding engine architecture.

These three problems are not merely technical inconveniences; they represent structural barriers that, if left unaddressed, confine LLM-based multi-agent simulation to the status of a toy demonstration rather than a reliable analytical tool.


\section{Application Context}

This simulation engine is positioned as a \textit{multi-stakeholder policy deliberation simulator}---a tool that enables policymakers to preview how diverse stakeholders might react to, argue against, and negotiate around proposed policies before committing to implementation.

The value proposition of this application is qualitatively distinct from quantitative policy modeling. Tools such as EUROMOD (European Commission Joint Research Centre) project fiscal impacts through microsimulation of tax-benefit systems, while platforms such as Polis (employed in Taiwan's vTaiwan initiative) aggregate real citizen opinions through asynchronous polling. The present engine occupies a different niche: it simulates the \textit{conversational dynamics} of a small group of identified stakeholders---their argumentative strategies, emotional responses, and coalition behaviors---in a controlled deliberation setting.

This positioning aligns naturally with the engine's demonstrated capabilities. The 20 benchmark scenarios are drawn from real-world policy and corporate contexts (EU GDPR enforcement, NYC Congestion Pricing, Boeing 737 MAX crisis response, among others), and the OCEAN personality model maps naturally to stakeholder communication styles: a high-Extraversion union leader advocates differently from a low-Extraversion compliance officer, even when both oppose the same policy.


\section{Contributions}

This study makes the following contributions:

\begin{enumerate}[leftmargin=*, itemsep=2pt]
    \item \textbf{Simulation engine architecture.} A six-component engine integrating a Persona State Controller (30 behavioral features, OCEAN trait estimation), a 4-Slot Candidate Pool (integrator, planner, challenger, skeptic), Sycophancy Detection, a Commitment Lifecycle tracker, an Action Layer (12 structured action types with world-state deltas), and Strategic Disposition calibration for multi-stakeholder policy deliberation.

    \item \textbf{Large-scale empirical validation.} A benchmark of 4,080 simulation runs across 20 real-world policy and corporate scenarios, three actor scales (3, 5, 10 agents), and two experimental conditions (engine-controlled vs.\ naive baseline), yielding statistically significant improvements in persona drift ($-5.7\%$, Cohen's $d = -0.94$), envelope violations ($-9.9\%$), and sycophancy reduction ($-69\%$).

    \item \textbf{Discovery of the RLHF behavioral ceiling.} Quantitative evidence that RLHF safety training imposes a fundamental limit on persona simulation fidelity, particularly for Neuroticism (\texttt{self\_doubt\_count} $= 0.000$ across both conditions) and low Agreeableness, establishing a boundary that cannot be overcome through engine-level intervention alone.

    \item \textbf{Iterative refinement methodology.} A follow-up 240-run iterative refinement demonstrating that benchmark-driven hyperparameter calibration (relationship delta clamping, fuzzy action validation) yields further improvements, establishing a reproducible refinement cycle for simulation engine development.
\end{enumerate}


\section{Report Structure}

The remainder of this report is organized as follows. Chapter~2 reviews related work in LLM-based persona simulation, multi-agent systems, the OCEAN personality model in NLP, behavioral fidelity evaluation, and existing policy simulation tools. Chapter~3 presents the system design of the simulation engine, detailing each of the six components with pseudocode for all deterministic logic. Chapter~4 describes the experimental setup, including scenarios, conditions, metrics, and statistical methods. Chapter~5 reports the results of the 4,080-run benchmark. Chapter~6 discusses the implications of these findings, the fidelity-authenticity tension, the RLHF ceiling, iterative refinement results, limitations, and future work. Chapter~7 concludes.
